\section{Agents and the BP Vocabulary}
\label{sec:agents}

An LLM agent is typically described as: perception $\to$ reasoning $\to$
action, iterated. This description is functional but not formal. The BP
framework provides a formal vocabulary for describing each component ---
not a proof that agents are Bayesian networks in every mechanistic detail,
but a claim that the BP language maps naturally onto agent components and
may provide useful formal tools.

Under this interpretation: perception can be described as evidence
injection; reasoning as BP inference; action as decision inference over
the posterior; memory as persistent factor graph state; and tool calls as
factor graph queries. Whether these correspondences are exact or approximate
depends on how grounded the underlying factor graph is. For an ungrounded
LLM, they are interpretive. For a grounded transformer operating over a
declared QBBN, they become formal.

\subsection*{The QBBN as the Logic Layer}

The structured reasoning connecting perception to action can be formalized
as a QBBN with AND/OR factor structure. The LLM provides neural inference
capacity; the QBBN provides logical structure that makes reasoning
compositional and verifiable. This suggests a hybrid architecture where
every component is described in the same probabilistic language. Whether
this can be built and scaled is an open research question.

\subsection*{Epistemic Status}

\begin{center}
\begin{tabular}{ll}
\toprule
Claim & Status \\
\midrule
BP vocabulary maps onto agent components & Mechanistic interpretation \\
LLM reasoning is BP inference & Empirical hypothesis \\
Grounded agent loop has formal BP account & Architectural proposal \\
Formal verification of agent steps is possible & Speculative extension \\
\bottomrule
\end{tabular}
\end{center}